\chapter{Relative coefficient envelopes}
\label{v4-ch:relative-coefficient-envelopes}

A coefficient algebra records which intermediate states may be
composed. Passing to a scalar trace can erase that information.
To retain it in a geometric realization, we shall attach the
coefficient algebra to a fixed geometric algebra and transport
bimodules together with their compositions.

The first calculation already determines the construction.
Let \(A\) and \(E\) be ordinary unital associative algebras over a field
\(k\), let \(P\) be a right \(E\)-module, and let \(Q\) be a left
\(E\)-module. Then
\[
 (A\otimes_kP)\otimes_{A\otimes_kE}(A\otimes_kQ)
 \longrightarrow A\otimes_k(P\otimes_EQ),
 \qquad
 (a\otimes p)\otimes(a'\otimes q)\longmapsto
 aa'\otimes(p\otimes q)
\]
is an isomorphism. Its inverse sends
\(a\otimes(p\otimes q)\) to
\((a\otimes p)\otimes(1_A\otimes q)\).
Balancing moves the second \(A\)-factor into the first without
interchanging them. Thus \(A\) need not be commutative.

For complexes, derived tensors retain the higher relations among
the coefficients. For geometric objects, the same calculation
uses the given geometric multiplication. We will construct the
comparison on the two-sided bars that define those derived
tensors. This specifies its signs, its unit, and its compatibility
with any finite string of interfaces.

\section{Coefficients in a monoidal category}
\label{v4-sec:relative-coefficient-setting}

Fix a commutative ring \(k\). Write \(D(k)\) for the stable
\(\infty\)-category of cochain complexes of \(k\)-modules
with quasi-isomorphisms inverted. Differentials have degree
\(1\); the shift satisfies
\(M[n]^i=M^{i+n}\).

Let \((\mathcal C,\star,K)\) be a stable presentable monoidal
\(\infty\)-category. Assume that \(\star\) preserves colimits separately in
both variables and that \(D(k)\) acts centrally on \(\mathcal C\).
Here centrality means that a scalar complex can pass through a
geometric factor, with the usual coherent Koszul interchange.
It supplies no interchange of two geometric factors.
We write \(X\otimes_kV\) for the action of \(V\in D(k)\) on
\(X\in\mathcal C\), and
\[
 \underline V=K\otimes_kV
\]
for the corresponding scalar object of \(\mathcal C\).
All tensor products in this chapter are derived unless an
ordinary representative is specified. Element formulas refer
to flat dg representatives of the structural maps.

Fix a unital associative algebra object \(A\) in \(\mathcal C\).
For a unital dg \(k\)-algebra \(B\), define
\begin{equation}\label{v4-eq:relative-algebra}
 R_B=A\otimes_kB.
\end{equation}
We call \(R_B\) the \emph{relative coefficient envelope}
with background \(A\).
The product uses central interchange followed by the two
multiplications. On homogeneous representatives it is
\[
 (a\otimes b)(a'\otimes b')
   =(-1)^{|b||a'|}aa'\otimes bb'.
\]
The algebra unit is the composite
\(K\to A\to A\otimes_kB\) obtained from the units of \(A\) and
\(B\). The object \(K\) remains the monoidal unit.

For another unital dg \(k\)-algebra \(E\),
a dg \(B\)--\(E\) bimodule \(P\) gives an
\(R_B\)--\(R_E\) bimodule
\begin{equation}\label{v4-eq:relative-interface}
 J_{B,E}(P)=A\otimes_kP
\end{equation}
with differential and actions
\begin{align}
 d(a\otimes p)
   &=d_Aa\otimes p+(-1)^{|a|}a\otimes d_Pp,
       \label{v4-eq:relative-differential}\\
 (a\otimes b)(x\otimes p)
   &=(-1)^{|b||x|}ax\otimes bp,
       \label{v4-eq:relative-left-action}\\
 (x\otimes p)(a'\otimes e)
   &=(-1)^{|p||a'|}xa'\otimes pe.
       \label{v4-eq:relative-right-action}
\end{align}
We call such a bimodule an \emph{interface} from its left
algebra to its right algebra. Two interfaces compose by
balanced tensor over their common algebra.

The signs in
\eqref{v4-eq:relative-left-action}--\eqref{v4-eq:relative-right-action}
are forced by the positions of the factors. For example, in a
right action the coefficient \(p\) passes through \(a'\), so an
odd \(p\) and an odd \(a'\) contribute a minus sign. A degree-\(r\)
coefficient bimodule map \(f:P\to P'\) is transported by
\begin{equation}\label{v4-eq:relative-map}
 J(f)(a\otimes p)=(-1)^{r|a|}a\otimes f(p).
\end{equation}
Associativity, the commuting left and right coefficient
actions, and the Koszul rule verify the bimodule and dg-functor
identities.

\section{The comparison of the two bars}
\label{v4-sec:relative-composition}

Let \(B,E,D\) be unital dg \(k\)-algebras, and let
\({}_BP_E\) and \({}_EQ_D\) be dg bimodules. The coefficient
derived tensor \(P\otimes_E^{\mathbb L}Q\) is the realization of
the simplicial two-sided bar
\begin{equation}\label{v4-eq:coefficient-bar}
 \mathcal B_n(P,E,Q)=P\otimes_kE^{\otimes_kn}\otimes_kQ.
\end{equation}
For \(n\geq1\), the faces are the right action on \(P\), the
\(n-1\) adjacent products in \(E\), and the left action on
\(Q\), in that order. Degeneracies insert \(1_E\).
In dg representatives we use the direct-sum total
\begin{equation}\label{v4-eq:bar-total}
 |\mathcal B_\bullet|
   =\bigoplus_{n\geq0}\mathcal B_n[n],
 \qquad
 d_{\mathrm{tot}}\big|_{\mathcal B_n[n]}
   =(-1)^nd_{\mathrm{int}}+\sum_{i=0}^{n}(-1)^id_i,
\end{equation}
where the face differential out of column \(0\) is zero.
The module structure on each shifted column is the shifted
module structure. In particular, a left action of degree
\(r\) acquires the sign \((-1)^{nr}\) in column \(n\).

Independently, the geometric balanced tensor
\[
 J_{B,E}(P)\star_{R_E}^{\mathbb L}J_{E,D}(Q)
\]
is the realization of
\begin{equation}\label{v4-eq:geometric-interface-bar}
 \mathcal L_n
   =(A\otimes_kP)\star
      (A\otimes_kE)^{\star n}\star(A\otimes_kQ).
\end{equation}
Its factors are in their displayed order. Define
\(\phi_n:\mathcal L_n\to A\otimes_k\mathcal B_n(P,E,Q)\) by
\begin{equation}\label{v4-eq:relative-bar-map}
\begin{split}
 &(a_0\otimes p)|(a_1\otimes e_1)|\cdots|
       (a_n\otimes e_n)|(a_{n+1}\otimes q)\\
 &\qquad\longmapsto
 (-1)^{\epsilon}(a_0a_1\cdots a_{n+1})
       \otimes(p|e_1|\cdots|e_n|q),\\
 &\epsilon
   =|p|\sum_{i=1}^{n+1}|a_i|
      +\sum_{j=1}^{n}|e_j|
          \sum_{i=j+1}^{n+1}|a_i|.
\end{split}
\end{equation}
Each summand in \(\epsilon\) records one scalar factor passing
through a later geometric factor. The \(A\)-factors are
multiplied in their original order.

\begin{lemma}\label{v4-lem:relative-bar-map}
The maps \(\phi_n\) form a simplicial dg bimodule map.
After realization they give a natural comparison
\[
 \Phi_{P,Q}:
 J_{B,E}(P)\star_{R_E}^{\mathbb L}J_{E,D}(Q)
 \longrightarrow
 J_{B,D}(P\otimes_E^{\mathbb L}Q).
\]
\end{lemma}

\begin{proof}
Consider an interior face, which combines the pairs
\((a_j,e_j)\) and \((a_{j+1},e_{j+1})\).
Multiplication before the shuffle contributes
\(|e_j||a_{j+1}|\). After multiplication these two entries
have degrees \(|a_j|+|a_{j+1}|\) and
\(|e_j|+|e_{j+1}|\). The difference between the two crossing
counts is exactly \(|e_j||a_{j+1}|\). Hence the face commutes
with \(\phi_n\). At the first and last faces the same count
gives the right and left module-action signs. Degeneracies
insert closed units of degree zero.

Central interchange and multiplication in \(A\) are chain
maps. They therefore commute with the internal differential.
The endpoint actions also commute with the map: both routes
move the same scalar factors through the same geometric
factors and then apply the same ordered product.

One totalization sign remains. The canonical chain map
\[
 |A\otimes_k\mathcal B_\bullet|
       \longrightarrow A\otimes_k|\mathcal B_\bullet|
\]
multiplies a term with geometric degree \(|a|\) in column
\(n\) by \((-1)^{n|a|}\).
For the differential of \(A\), changing \(|a|\) by one
contributes \((-1)^n\), as required by
\eqref{v4-eq:bar-total}. For a face, changing \(n\) to \(n-1\)
contributes \((-1)^{|a|}\), the tensor-differential sign on
the target. The internal coefficient differential is
unchanged. Composing this identification with
\(|\phi_\bullet|\) gives the asserted dg bimodule map.
\end{proof}

\begin{theorem}[Relative coefficient composition]
\label{v4-thm:relative-composition}
Under the hypotheses of
Section~\ref{v4-sec:relative-coefficient-setting},
the comparison \(\Phi_{P,Q}\) is an equivalence.
For every finite string of coefficient bimodules, the
comparisons satisfy the associativity and unit coherences.
The identity coefficient bimodule \(B\) is sent to the
regular interface \(R_B\).
\end{theorem}

\begin{proof}
Resolve \(Q\) by the augmented free left-\(E\) bar
\[
 F_n=E\otimes_kE^{\otimes_kn}\otimes_kQ
       \longrightarrow Q.
\]
After forgetting its left \(E\)-action, insertion of \(1_E\)
at the left gives an extra degeneracy. The augmented
simplicial identities contract its realization, so
\(|F_\bullet|\simeq Q\).

Tensoring with \(A\) gives a resolution of \(J_{E,D}(Q)\).
Its term in degree \(n\) is the free left-\(R_E\) module
\[
 R_E\star\underline{E^{\otimes_kn}\otimes_kQ}.
\]
Balancing it with \(J_{B,E}(P)\) leaves
\[
 A\otimes_kP\otimes_kE^{\otimes_kn}\otimes_kQ.
\]
The resulting faces are exactly the faces of
\eqref{v4-eq:coefficient-bar}. Their realization is
\(J_{B,D}(P\otimes_E^{\mathbb L}Q)\).

To identify the particular comparison, map the full free
\(R_E\)-bar resolution of \(J(Q)\) to \(A\otimes_kF_\bullet\)
by multiplying its \(A\)-entries in order.
The crossing signs are those in
\eqref{v4-eq:relative-bar-map}. The preceding lemma verifies
the faces, degeneracies, differentials, and outer actions.
Both augmentations are equivalences; hence the realized
comparison between the resolutions is an equivalence.
Derived balancing preserves this equivalence. The resulting
map is \(\Phi_{P,Q}\).

For three interfaces, use their common bisimplicial bar.
Both parenthesizations move each scalar coefficient through
every later geometric entry and then multiply the geometric
entries in order. Already in bar degree zero the exponents
are
\[
 |p||a'|+(|p|+|q|)|a''|
 \quad\text{and}\quad
 |q||a''|+|p|(|a'|+|a''|),
\]
which agree. In arbitrary bar degrees the two exponents count
the same pairs of entries.

The same construction on the common multisimplicial bar of
any finite string gives one comparison before a
parenthesization is chosen. Every parenthesization induces
that comparison, so all associativity diagrams commute.
For units, \(J_{B,B}(B)=R_B\) with its regular actions, and
the two unit comparisons are the augmentations contracted
by insertion of \(1_B\). These are the same degeneracies
in every multibar. Realization therefore supplies the unit
and associativity coherences together.
\end{proof}

The proof uses the geometric product of \(A\), scalar
interchange, and colimits. Consequently, any additional
operation that commutes with these maps commutes with
\(\Phi_{P,Q}\). In a collision model this requires the
specified collision comparisons for \(A\); the algebraic
formula alone does not construct them.

\section{Recovery at the empty configuration}
\label{v4-sec:relative-empty-recovery}

Suppose that \(\mathcal C\) has a colimit-preserving strong
monoidal functor
\[
 e_0:\mathcal C\longrightarrow D(k)
\]
compatible with scalar action, and fix an algebra
equivalence \(e_0(A)\simeq k\). In a word-convolution
category, \(e_0\) is evaluation at the empty configuration.

\begin{proposition}\label{v4-prop:relative-retraction}
There are natural identifications
\[
 e_0(R_B)\simeq B,\qquad e_0(J_{B,E}(P))\simeq P.
\]
Under these identifications,
\(e_0(\Phi_{P,Q})\) is the identity comparison on
\(P\otimes_E^{\mathbb L}Q\).
Empty evaluation also recovers every transported coefficient
map and every diagram of such maps.
\end{proposition}

\begin{proof}
Compatibility with scalar action gives
\(e_0(A\otimes_kP)\simeq e_0(A)\otimes_kP\simeq P\).
Since \(e_0\) is strong monoidal and preserves colimits,
it carries the geometric bar to the coefficient bar.
The product of the \(A\)-entries in
\eqref{v4-eq:relative-bar-map} becomes multiplication of
copies of the scalar unit. The comparison is therefore
the identity under the canonical unit identifications.
The same argument applied to
\eqref{v4-eq:relative-map} recovers coefficient maps and
their compositions.
\end{proof}

In particular, a homotopy between two transported coefficient
maps gives a homotopy between the original maps by applying
\(e_0\). Recovery does
not assert that every geometric bimodule map is induced
by a coefficient map. The geometric category can contain
additional maps.

The empty algebra of \(R_B\) is \(B\). This retained
coefficient is the reason that arbitrary coefficient
interfaces compose with their original units. Replacing
it by \(k\) changes the comparison problem.

\section{Endomorphisms and the orientation of composition}
\label{v4-sec:endomorphism-interfaces}

A useful coefficient system is supplied by endomorphisms
of one algebra. Let \(B\) be an ordinary unital
\(k\)-algebra, and let \(\alpha:B\to B\) be a unital
endomorphism. Define
\[
 Q_\alpha={}_\alpha B_1,
 \qquad b\cdot x\cdot c=\alpha(b)xc.
\]
For endomorphisms \(\alpha,\beta\), an element \(t\in B\)
satisfying
\[
 t\alpha(b)=\beta(b)t\qquad(b\in B)
\]
defines a bimodule map \(Q_\alpha\to Q_\beta\),
\(x\mapsto tx\). Thus actual intertwiners, including their
relations, become maps of interfaces.

\begin{proposition}\label{v4-prop:endomorphism-composition}
There is a natural bimodule isomorphism
\[
 Q_\alpha\otimes_BQ_\beta
       \xrightarrow{\ \sim\ }Q_{\beta\circ\alpha},
 \qquad x\otimes y\longmapsto\beta(x)y.
\]
It also computes the derived tensor. Consequently,
\[
 J(Q_\alpha)\star_{R_B}^{\mathbb L}J(Q_\beta)
       \simeq J(Q_{\beta\circ\alpha})
\]
with the coherences of
Theorem~\ref{v4-thm:relative-composition}.
\end{proposition}

\begin{proof}
The balancing relation is
\(xb\otimes y=x\otimes\beta(b)y\).
Its two images are both \(\beta(x)\beta(b)y\).
The left action of \(b\) on the image is
\(\beta(\alpha(b))\), and the right action is unchanged.
The inverse is \(y\mapsto1_B\otimes y\).
The right \(B\)-module underlying \(Q_\alpha\) is free
of rank one, so its ordinary tensor computes the
derived tensor. The last assertion is the relative
composition theorem.
\end{proof}

Composition is therefore \(\beta\circ\alpha\) in the
displayed order. A tensor category of endomorphisms using
\(\alpha\circ\beta\) must have its tensor order reversed
before this functor is called monoidal. This convention
also determines how associators and duality maps are
transported.

Every specified identity among intertwiner diagrams
survives under \(J\), including pentagon and evaluation
identities when they occur in the input. Neither a
braiding nor a tensor product of independent
representations is supplied by balanced interface
composition.

\section{Scalar change}
\label{v4-sec:relative-scalar-change}

Let \(k\to k'\) be a map of commutative rings.
Assume that a colimit-preserving monoidal scalar-change
functor from \(\mathcal C\) to a corresponding category
\(\mathcal C'\) is given and is compatible with the
central scalar actions. Denote it by \(X\mapsto X_{k'}\).
Then
\[
 (R_B)_{k'}
   \simeq A_{k'}\otimes_{k'}
                (B\otimes_k^{\mathbb L}k'),
 \qquad
 J(P)_{k'}
   \simeq A_{k'}\otimes_{k'}
                (P\otimes_k^{\mathbb L}k').
\]
These identifications preserve \(\Phi_{P,Q}\).
Indeed, scalar change preserves the two bar realizations
and their common crossing map. Its coherence with a
second scalar change follows on the same diagrams.

For arithmetic coefficient towers, this supplies each
finite-precision comparison once the geometric
scalar-change functors have been constructed.
Continuity of an inverse limit is a further statement:
it requires the chosen completed tensor and the
corresponding limit comparison. The finite bar formulas
do not interchange an arbitrary limit with realization.
